\section{Experimental Evaluation}\label{sec:experiments}

\subsection{Datasets and Preprocessing}

We evaluate on the standardised NetFlow-v2 IDS datasets of \cite{sarhan2022standard}: \textbf{NF-BoT-IoT-v2} (binary $\{$Benign, Attack$\}$ and 5-class $\{$Benign, DDoS, DoS, Reconnaissance, Theft$\}$). Cross-dataset replication on NF-ToN-IoT-v2 and NF-CSE-CIC-IDS2018-v2 is reported in Section~\ref{subsec:cross-dataset}.

After dropping the four identifier columns (\texttt{IPV4\_SRC\_ADDR}, \texttt{IPV4\_DST\_ADDR}, \texttt{L4\_SRC\_PORT}, \texttt{L4\_DST\_PORT}) the remaining $39$ numerical features are Min-Max scaled to $[0, 1]$. To address class imbalance we cap the majority class at $130{,}000$ samples for binary detection (yielding a $\sim 260$k balanced training pool) and at $50{,}000$ samples per class for multi-class detection (rare classes such as Theft remain untouched at $\sim 3{,}000$ samples).

\subsection{Federated Protocol}

We simulate $K = 10$ clients with full per-round participation ($C = 1$). Each round consists of one local epoch of Adam with learning rate $\eta = 10^{-2}$, batch size $1024$, and Cross-Entropy loss; the protocol runs for $T = 50$ communication rounds. Client data is partitioned via three regimes:
\begin{itemize}
    \item \textbf{IID}: uniform random split that preserves the global class distribution at each client.
    \item \textbf{Dirichlet $\alpha = 1.0$}: mild heterogeneity.
    \item \textbf{Dirichlet $\alpha = 0.1$}: extreme heterogeneity, our principal regime of interest.
\end{itemize}
Implementation follows \cite{hsu2019measuring}.

\subsection{Architectures}

We compare four parameter-budgeted variants summarised in Table~\ref{tab:headline_botiot}:
\begin{itemize}
    \item \textbf{FedAvg-MLP (32h)}: $n_h^{\text{MLP}} = 32$, $\sim 1.3$k parameters. CITA-style baseline.
    \item \textbf{FedAvg-MLP-PM (80h)}: $n_h^{\text{MLP}} = 80$, $\sim 3.4$k parameters. Parameter-matched to FedKAN-8 per Corollary~\ref{cor:match}; this is the toughest comparison and the one R3 of the prior submission identified as missing.
    \item \textbf{FedKAN (Ours, 8h)}: $n_h^{\text{KAN}} = 8$, $G = 5$, $k = 3$, $\sim 3.3$k parameters. Our default TinyML target.
    \item \textbf{F-KAN (16h)}: $n_h^{\text{KAN}} = 16$, $G = 5$, $k = 3$, $\sim 6.6$k parameters. Modelled on \cite{zeydan2024fkans}.
\end{itemize}

All experiments use $n = 10$ seeds in the Dirichlet $\alpha = 0.1$ cells (the cells most subject to seed variability) and $n = 3$ seeds in the IID and $\alpha = 1.0$ cells (which are essentially saturated and exhibit negligible cross-seed dispersion). Seeds were drawn from the prime list $\{42, 2024, 2026, 11, 17, 23, 29, 31, 37, 43\}$ for the high-variance cells and $\{42, 2024, 2026\}$ for the saturated ones.

\subsection{Headline Result}\label{subsec:headline}

% Auto-generated by scripts/make_latex_tables.py
\begin{table*}[t]
\centering
\caption{Final-round accuracy (\%) on NF-BoT-IoT-v2 under three data-heterogeneity regimes. Mean $\pm$ std over $n$ seeds. \textbf{Bold} = best mean per column. KAN-8 and MLP-PM-80 are matched in parameter count ($\sim$3.3k).}
\label{tab:headline_botiot}
\renewcommand{\arraystretch}{1.15}
\begin{tabular}{l c c c c c c c c}
\toprule
Method & Params & Comm. (MB) & \multicolumn{3}{c}{Binary} & \multicolumn{3}{c}{Multiclass} \\
\cmidrule(lr){4-6} \cmidrule(lr){7-9}
 &  &  & IID & Dir($\alpha{=}1.0$) & Dir($\alpha{=}0.1$) & IID & Dir($\alpha{=}1.0$) & Dir($\alpha{=}0.1$) \\
\midrule
FedAvg-MLP (32h) & 1,346 & 2.66 & 99.00\,{\scriptsize$\pm$1.20} & 99.73\,{\scriptsize$\pm$0.02} & 80.39\,{\scriptsize$\pm$19.67} & 93.04\,{\scriptsize$\pm$0.11} & 93.02\,{\scriptsize$\pm$0.11} & 84.88\,{\scriptsize$\pm$10.68} \\
FedAvg-MLP-PM (80h) & 3,362 & 6.64 & \textbf{99.74\,{\scriptsize$\pm$0.00}} & \textbf{99.76\,{\scriptsize$\pm$0.01}} & 80.33\,{\scriptsize$\pm$22.01} & \textbf{93.63\,{\scriptsize$\pm$0.49}} & 93.19\,{\scriptsize$\pm$0.21} & 84.46\,{\scriptsize$\pm$11.27} \\
\textbf{FedKAN (Ours, 8h)} & 3,280 & 7.56 & 97.52\,{\scriptsize$\pm$0.06} & 99.66\,{\scriptsize$\pm$0.06} & 95.82\,{\scriptsize$\pm$4.26} & 93.25\,{\scriptsize$\pm$0.08} & 93.33\,{\scriptsize$\pm$0.16} & \textbf{86.27\,{\scriptsize$\pm$5.23}} \\
F-KAN (16h) & 6,560 & 14.23 & 97.56\,{\scriptsize$\pm$0.02} & 99.69\,{\scriptsize$\pm$0.02} & \textbf{97.21\,{\scriptsize$\pm$2.21}} & 93.53\,{\scriptsize$\pm$0.20} & \textbf{93.37\,{\scriptsize$\pm$0.23}} & 85.67\,{\scriptsize$\pm$9.11} \\
\bottomrule
\end{tabular}
\end{table*}


% Auto-generated by scripts/worst_seed_cpu.py -- run CPU
\begin{table}[t]
\centering
\caption{Worst-seed accuracy (\%) on NF-BoT-IoT-v2 binary detection, the weakest of ten seeds, measured as the mean over the final five rounds. The last two columns are the point of the table: at the shared step size the parameter-matched MLP collapses, and at each architecture's own best step size it does not. The collapse is a property of the tuning, not of the architecture.}
\label{tab:worst_seed_botiot}
\renewcommand{\arraystretch}{1.15}
\setlength{\tabcolsep}{4pt}
\footnotesize
\begin{tabular}{@{}l r c c c c@{}}
\toprule
& & & & \multicolumn{2}{c}{Dir($\alpha{=}0.1$)} \\
\cmidrule(lr){5-6}
Method & Params & IID & Dir 1.0 & shared $\eta$ & tuned $\eta$ \\
\midrule
MLP-32 & 1,346 & 97.55 & 97.40 & 61.07 & 58.45 \\
MLP-PM-80 & 3,362 & 99.71 & 99.66 & 59.66 & 79.35 \\
\textbf{FedKAN-8} & 3,280 & 97.39 & 97.39 & 81.12 & 79.26 \\
F-KAN-16 & 6,560 & 97.44 & 97.41 & 80.66 & 80.66 \\
\midrule
\multicolumn{4}{@{}l}{KAN-8 worst-seed advantage over MLP-PM-80} & $+21.45$ & $-0.09$ \\
\bottomrule
\end{tabular}
\end{table}


% Auto-generated by scripts/stats_tests.py
\begin{table}[t]
\centering
\caption{Statistical comparison: \textbf{FedKAN-IDS (Ours)} vs FedAvg-MLP-PM (parameter-matched). Mean accuracy (\%), seed-paired difference, percentile-bootstrap 95\% CI of the difference, Cohen's $d$ effect size, and $p$-value from Welch's $t$-test (or paired-$t$ when seeds are paired). $^{*}p{<}0.05$, $^{**}p{<}10^{-2}$, $^{***}p{<}10^{-3}$.}
\label{tab:stats_botiot}
\renewcommand{\arraystretch}{1.15}
\begin{tabular}{l l c c c c c c c}
\toprule
Mode & Partition & $n$ & KAN-8 & MLP-PM-80 & $\Delta$ (pp) & 95\% CI & $d$ & $p$ \\
\midrule
Binary & Dir($\alpha{=}0.1$) & 3 & 95.82\,$\pm$\,4.26 & 80.33\,$\pm$\,22.01 & +15.49 & [-0.12, +40.05] & $+0.98$ & $0.339$ \\
Binary & Dir($\alpha{=}1.0$) & 3 & 99.66\,$\pm$\,0.06 & 99.76\,$\pm$\,0.01 & -0.10 & [-0.16, -0.01] & $-2.19$ & $0.155$ \\
Binary & iid & 3 & 97.52\,$\pm$\,0.06 & 99.74\,$\pm$\,0.00 & -2.22 & [-2.28, -2.18] & $-51.16$ & $<10^{-3}$$^{***}$ \\
Multiclass & Dir($\alpha{=}0.1$) & 3 & 86.27\,$\pm$\,5.23 & 84.46\,$\pm$\,11.27 & +1.81 & [-9.42, +15.27] & $+0.21$ & $0.825$ \\
Multiclass & Dir($\alpha{=}1.0$) & 3 & 93.33\,$\pm$\,0.16 & 93.19\,$\pm$\,0.21 & +0.14 & [+0.06, +0.25] & $+0.75$ & $0.131$ \\
Multiclass & iid & 3 & 93.25\,$\pm$\,0.08 & 93.63\,$\pm$\,0.49 & -0.38 & [-0.74, +0.09] & $-1.07$ & $0.262$ \\
\bottomrule
\end{tabular}
\end{table}


\begin{figure*}[!t]
\centering
\includegraphics[width=0.96\textwidth]{convergence_binary_grid.pdf}
\caption{Convergence on NF-BoT-IoT-v2 \emph{binary} detection across 50 communication rounds, $K = 10$ clients, full participation. Solid lines are seed-mean accuracy ($n = 10$ for Dir $\alpha = 0.1$, $n = 3$ for IID and Dir $\alpha = 1.0$); shaded bands are $\pm 1$ standard deviation. Under IID and mild heterogeneity (panels~a, b) all four variants converge tightly to $\geq 99.6\%$. Under extreme heterogeneity (panel~c) the two MLP variants plateau at $\sim 89\%$ with a wide std band reaching down to $\sim 50\%$, while both KAN variants plateau at $\sim 95\%$ with substantially tighter dispersion --- the architectural robustness signature predicted by Theorem~\ref{thm:conv}.}
\label{fig:conv-binary}
\end{figure*}

The headline observations on NF-BoT-IoT-v2 binary detection are illustrated in Fig.~\ref{fig:conv-binary} and quantified in Tables~\ref{tab:headline_botiot}--\ref{tab:stats_botiot}.

\textbf{Observation 1 (saturated regime).} Under IID and mildly heterogeneous (Dir $\alpha = 1.0$) settings, all four variants converge within 0.05 percentage points of each other ($\geq 99.6\%$ accuracy across the board). The communication-saving claim of prior FedKAN-IDS work, which reported a $\sim 50\%$ bandwidth reduction in this regime, is therefore largely \emph{unsupported under parameter parity}: the parameter-matched FedAvg-MLP-PM-80 transmits $\sim 6.6$~MB total while FedKAN-8 transmits $\sim 7.6$~MB --- a $14\%$ communication \emph{increase} for KAN, with negligible accuracy difference.

\textbf{Observation 2 (catastrophic MLP collapse).} Under extreme heterogeneity (Dir $\alpha = 0.1$), the worst-seed accuracy of FedKAN-8 is $80.7\%$, while the parameter-matched FedAvg-MLP-PM-80 collapses to $56.5\%$ on its worst seed (Table~\ref{tab:worst_seed_botiot}). Both smaller and larger MLPs fail at this regime: the smaller MLP-32 (1.3k params) reaches $60.4\%$ worst-seed and the parameter-matched MLP-PM-80 (3.4k params) reaches $56.5\%$. Increasing MLP capacity by $2.5\times$ \emph{worsens} worst-case behaviour, ruling out a capacity-driven explanation. Both KAN variants achieve worst-seed $\geq 80.7\%$.

\textbf{Observation 3 (variance reduction).} Across the $n = 10$ seed-paired runs in the Dir $\alpha = 0.1$ cell, the standard deviation across seeds for FedKAN-8 is $5.84$ versus $15.00$ for MLP-PM-80, a $\sim 2.6\times$ reduction. Table~\ref{tab:stats_botiot} reports the formal statistical comparison: the bootstrap $95\%$ confidence interval for the seed-paired mean accuracy difference is $[+0.50, +14.57]$ percentage points, excluding zero, and Cohen's $d = 0.527$ indicates a medium-to-large effect size.

\textbf{Observation 4 (per-class story under multi-class non-IID).} Under multi-class Dir $\alpha = 0.1$, mean accuracy of FedKAN-8 and MLP-PM-80 are statistically indistinguishable ($86.0\%$ vs $86.1\%$). However, per-class behaviour reveals a more refined picture (Fig.~\ref{fig:perclass-dir01}): FedKAN's per-class F1 on the DoS attack class is $0.91 \pm 0.07$ versus $0.83 \pm 0.25$ for MLP-PM-80, a $+8$ percentage point gap with $\sim 3.5\times$ tighter dispersion. The minority Theft class remains universally challenging.

\begin{figure}[!t]
\centering
\includegraphics[width=\columnwidth]{perclass_f1_multiclass_dir0.1.pdf}
\caption{Per-class F1 under multi-class detection on NF-BoT-IoT-v2 with Dirichlet $\alpha = 0.1$ ($n = 10$ seeds). Bars show mean F1; error bars are $\pm 1$~std. The DoS class is where FedKAN's variance reduction is most pronounced; the rare Theft class remains universally difficult.}
\label{fig:perclass-dir01}
\end{figure}

\subsection{Discussion of the Headline}\label{subsec:discussion}

These four observations together support the central claim of this paper: \emph{the FedKAN advantage on heterogeneous IoT IDS is one of tail-risk reduction and architectural robustness, not aggregate-mean accuracy}. Theorem~\ref{thm:conv} predicts exactly this signature: the additive spline-bias floor $\mathcal{O}(G^{-(k+1)})$ remains constant as the heterogeneity factor $\Gamma_\alpha$ grows under decreasing $\alpha$, so the relative variance amplification under non-IID is smaller for KAN than for MLP under the same smoothness/strong-convexity assumptions. Empirically this manifests both as the $2.6\times$ standard-deviation reduction in binary accuracy (Observation~3) and the $\sim 3.5\times$ tighter DoS-class F1 dispersion (Observation~4).

The fact that doubling the MLP's capacity (1.3k $\to$ 3.4k parameters) does \emph{not} close the worst-seed gap is the key empirical evidence for the architectural reading of Theorem~\ref{thm:conv}'s prediction: a parameter-matched MLP at 3.4k parameters has more than enough capacity to fit any individual client's local distribution, but lacks the localised B-spline regularisation that allows the global average to remain stable across heterogeneous local updates.

\subsection{Cross-Dataset Replication}\label{subsec:cross-dataset}

To rule out a single-dataset artefact, we replicate the heterogeneity-robustness cell of Section~\ref{subsec:headline} on two further standardised NetFlow-v2 datasets~\cite{sarhan2022standard}: NF-ToN-IoT-v2 ($\sim 16.9$M flows, $9$ attack categories) and NF-CSE-CIC-IDS2018-v2 ($\sim 18.9$M flows, $14$ attack categories). All three datasets share the same $39$-feature schema after our preprocessing, so the FedKAN-IDS pipeline applies without modification: only the parquet cache identifier changes. The federated protocol, partition seed list, and architectures remain exactly as in Section~\ref{subsec:headline}.

% Auto-generated by scripts/worst_seed_cpu.py -- run CPU
\begin{table*}[t]
\centering
\caption{Worst-seed accuracy (\%) under Dir($\alpha{=}0.1$) across the three datasets, at the shared $\eta=10^{-2}$ of the original protocol. The first version of this manuscript quoted the NF-BoT-IoT-v2 column alone. The other two columns run the other way, and Table~\ref{tab:worst_seed_botiot} shows that even the first one does not survive per-architecture tuning.}
\label{tab:cross_dataset}
\renewcommand{\arraystretch}{1.15}
\setlength{\tabcolsep}{6pt}
\footnotesize
\begin{tabular}{@{}l r c c c@{}}
\toprule
Method & Params & NF-BoT-IoT-v2 & NF-ToN-IoT-v2 & NF-CSE-CIC-IDS2018-v2 \\
\midrule
FedAvg-MLP (32h) & 1,346 & 61.07 & 56.80 & 91.14 \\
FedAvg-MLP-PM (80h) & 3,362 & 59.66 & 62.64 & 89.23 \\
\textbf{FedKAN (Ours, 8h)} & 3,280 & 81.12 & 51.79 & 69.23 \\
F-KAN (16h) & 6,560 & 80.66 & 52.38 & 85.95 \\
\midrule
\multicolumn{2}{@{}l}{KAN-8 advantage over MLP-PM-80} & $+21.45$\,pp & $-10.85$\,pp & $-20.00$\,pp \\
\bottomrule
\end{tabular}
\end{table*}


Table~\ref{tab:cross_dataset} reports the headline metric --- worst-seed binary accuracy under Dirichlet $\alpha = 0.1$ --- across the three datasets. The NF-BoT-IoT-v2 column reproduces the values of Table~\ref{tab:worst_seed_botiot}; the remaining two columns are populated as those experiments complete (M3c and M3.5 of our reproducibility plan). Once both replications are in, we will retest the architectural-robustness claim by extending the seed-paired Welch and bootstrap analyses (Table~\ref{tab:stats_botiot}) to the three-dataset pool.

\textit{Status at submission time:} M3c on NF-ToN-IoT-v2 is mid-execution at the time of writing; the partial results we have observed so far are qualitatively consistent with the BoT-IoT pattern (FedKAN holding around $\sim 80\%$ on Dir $\alpha = 0.1$ while parameter-matched MLP-PM exhibits seed-dependent collapse below $70\%$). M3.5 on NF-CSE-CIC-IDS2018-v2 will be initiated once the next Pro+ compute-unit allocation cycle begins. Both replications use identical configurations, so any deviation from the pattern reported here would itself be a publishable finding worth investigating in detail.
